\chapter{Stand der Forschung und Forschungslücke}

\section{Task-adaptive Repräsentationen und Label-Effizienz}

Ein Forschungsstrang der Phishing-Erkennung untersucht, wie transformerbasierte
Repräsentationen bereits vor der gelabelten Zielklassifikation an Web- und insbesondere
URL-Daten angepasst werden können. URLTran vergleicht zwei Ansätze: ausschließlich
auf URLs vortrainierte und anschließend feinabgestimmte Transformer sowie die direkte
Feinabstimmung öffentlich vortrainierter BERT- und RoBERTa-Modelle. Die Evaluation
berücksichtigt gezielt Betriebspunkte mit sehr niedrigen False-Positive-Raten
(vgl. [Maneriker et al. (2021), PDF-S. 1--2]). PhishBERT wird auf mehr als drei Milliarden
ungelabelten URLs vortrainiert und anschließend für die Phishing-Klassifikation
weiterverwendet (vgl. [Wang et al. (2023), Abstract, o.\,S.]). Auch urlBERT nutzt mehrere Milliarden
URLs und kombiniert mehrere selbstüberwachte Vortrainingsziele; ergänzend wird die
Sensitivität gegenüber reduziertem gelabeltem Trainingsumfang untersucht
(vgl. [Li et al. (2025), S. 1--2 und 12]).

Neben der domänennäheren Repräsentationsbildung rückt damit die verfügbare gelabelte
Downstream-Supervision in den Fokus. Gui et al. führen den Aufwand der Erhebung und
Annotation großer gelabelter Datenbestände als eine Motivation für Self-Supervised Learning
an (vgl. [Gui et al. (2024), S. 1]). Für spezialisierte URL-Repräsentationen untersuchen
Li et al. ergänzend, wie sich die Downstream-Leistung bei reduziertem gelabeltem
Trainingsumfang verändert (vgl. [Li et al. (2025), S. 12]).

Für die vorliegende Arbeit sind daraus zwei Aspekte relevant. Erstens wird untersucht,
welchen zusätzlichen Nutzen die in Abschnitt 2.3.3 eingeführte task-adaptive
Weiteranpassung bereits vortrainierter URL- und HTML-Text-Encoder für die
Phishing-Klassifikation erzeugt. Durch die getrennte und gemeinsame Anpassung beider
Informationszweige wird zugleich betrachtet, wie sich dieser Effekt zwischen URL und
HTML-Text verteilt. Zweitens wird die gelabelte Downstream-Supervision systematisch
über mehrere Labelbudgets variiert, um zu bestimmen, wie sich der
Repräsentationsvorteil mit wachsendem Trainingsumfang entwickelt und ab welchem
reduzierten Labelbudget das festgelegte Leistungsniveau vollständig trainierter
Referenzen erreicht wird.

\section{Informationsgrundlage und Systemarchitektur}

Neben der Repräsentationsbildung unterscheiden sich bestehende Verfahren hinsichtlich
der für die Klassifikation verwendeten Webseiteninformationen. Aljofey et al. kombinieren
URL-Zeichenfolgen, Hyperlinkinformationen und textuelle HTML-Merkmale und führen diese
einem XGBoost-Klassifikator zu
(vgl. [Aljofey et al. (2022), S. 1--2 und 13]). Opara et al. verarbeiten URL- und
HTML-Informationen dagegen in einem gemeinsamen neuronalen Verfahren
(vgl. [Opara et al. (2024), S. 2]). Damit werden URL und HTML-Inhalt als unterschiedliche
Informationsquellen derselben Webseite für die Klassifikation nutzbar gemacht.

Über den textuellen HTML-Inhalt hinaus kann auch die Dokumentstruktur berücksichtigt
werden. Yoon et al. kombinieren URL-Repräsentationen mit einem graphbasierten
HTML-DOM und untersuchen den zusätzlichen Beitrag der strukturellen Information anhand
von Ablationen (vgl. [Yoon et al. (2024), S. 2, 13 und 15]). Die bestehende Forschung
liefert damit eine Grundlage dafür, adressbezogene, textuelle und strukturelle
Informationen innerhalb eines gemeinsamen Klassifikationssystems zu berücksichtigen.

Die vorliegende Arbeit trennt dabei Repräsentations- und Systemebene. Zunächst wird die
task-adaptive URL- und HTML-Text-Repräsentation unter konstant gehaltenen
Downstream-Bedingungen untersucht. Die daraus ausgewählte Repräsentationsbasis wird
anschließend schrittweise um weitere Systemkomponenten erweitert. Eine sequenzielle
Ablation von URL über HTML-Text und DOM bis zur lernbaren Gating-Fusion bestimmt
hierbei den inkrementellen Beitrag der zusätzlichen Informations- und
Fusionskomponenten.

Neben dem Informationsumfang kann auch die Verarbeitungstiefe Bestandteil des
Systementwurfs sein. Liu et al. beschreiben ein mehrstufiges
Phishing-Erkennungsverfahren und berichten eine Verkürzung der Ausführungszeit
bei zugleich verbesserter Erkennung
(vgl. [Liu et al. (2021), Abstract, o.\,S.]). Dieses Prinzip bildet den konzeptionellen
Bezug für die URL-first-Kaskade der vorliegenden Arbeit. Die URL-Stufe dient dabei
als vorgelagerter Entscheidungspfad, während für unsichere Fälle weiterhin das
vollständige multimodale System ausgeführt wird.

\section{Generalisierung und betriebliche Evaluationsbedingungen}

Die Aussagekraft eines Phishing-Klassifikators hängt neben der Modellarchitektur auch
von der Zusammensetzung der Evaluationsdaten ab. Rashid et al. untersuchen
URL-Klassifikatoren über unterschiedliche Datensätze hinweg und berichten deutliche
Leistungsunterschiede zwischen datensatzinternen und datensatzübergreifenden
Evaluationen (vgl. [Rashid et al. (2024), Abstract, o.\,S.]). Dalton et al. adressieren im
PhreshPhish-Benchmark eine verwandte Problematik durch zeitlich getrennte Trainings-
und Testdaten sowie eine zusätzliche Reduktion stark ähnlicher Instanzen
(vgl. [Dalton et al. (2026), S. 5--6]).

Für die vorliegende Untersuchung wird die Leistungsrelation zwischen den
Repräsentationsvarianten und Labelbudgets deshalb nicht ausschließlich auf dem
vollständigen Testbestand betrachtet. Ergänzende domain-, template- und zeitbezogene
Teilmengen bilden zusätzliche Evaluationsbedingungen. Sie dienen der Prüfung, ob die
zuvor beobachteten Repräsentations- und Budgetrelationen auch unter veränderter
Zusammensetzung der Testdaten bestehen bleiben.

Eine weitere Evaluationsdimension bildet der Klassifikationsbetriebspunkt. URLTran
untersucht seine Modelle gezielt bei niedrigen False-Positive-Raten
(vgl. [Maneriker et al. (2021), PDF-S. 1--2]). Auch im operativen Sicherheitsbetrieb
werden False-Positive-Raten gemeinsam mit Größen wie Eskalationen und
Bearbeitungsaufwand als Qualitätskennzahlen betrachtet
(vgl. [ENISA (2020), S. 7]). Die vorliegende Arbeit verwendet daher einen auf einer
getrennten Kalibrationsmenge bestimmten Low-FPR-Betriebspunkt, der anschließend
unverändert auf die finalen Evaluationsbedingungen übertragen wird.

Für die technische Bewertung des Gesamtsystems kommen Ressourcenkennzahlen der
Inferenz hinzu. Reddi et al. unterscheiden insbesondere latenz- und
durchsatzorientierte Inferenzszenarien
(vgl. [Reddi et al. (2020), S. 5--6]). Entsprechend werden neben der
Klassifikationsleistung auch Latenz, Durchsatz, GPU-Speicherbedarf und
Eskalationsrate betrachtet.

\section{Forschungslücke und Einordnung der Arbeit}

Die betrachteten Forschungsarbeiten adressieren die für diese Arbeit relevanten
Aspekte überwiegend auf einzelnen Ebenen. Spezialisierte selbstüberwachte
URL-Repräsentationen zeigen die Nutzbarkeit domänennäherer Repräsentationsbildung;
Untersuchungen reduzierter Labelmengen betrachten die Downstream-Leistung bei
begrenzter Supervision. Weitere Arbeiten kombinieren mehrere Webseiteninformationen
oder verwenden mehrstufige Klassifikationsarchitekturen.

Aus dieser Ausgangslage ergibt sich für die vorliegende Arbeit die Verbindung dreier
Untersuchungsebenen. Zunächst wird der zusätzliche Effekt einer task-adaptiven
Weiteranpassung bereits vortrainierter URL- und HTML-Text-Repräsentationen isoliert.
Darauf aufbauend wird untersucht, wie sich dieser Effekt über unterschiedliche
Labelbudgets entwickelt und ab welchem reduzierten Budget das festgelegte
Leistungsniveau vollständig trainierter Referenzen erreicht wird. Die ausgewählte
Repräsentationsbasis wird anschließend in ein multimodales Gesamtsystem überführt,
dessen zusätzliche Komponenten und gestufte Verarbeitung getrennt evaluiert werden.

Tabelle~\ref{tab:forschungsableitung} fasst die für das Untersuchungsdesign
maßgeblichen Forschungsbezüge zusammen.

\begin{table}[!ht]
\centering
\caption{Einordnung des Untersuchungsdesigns in den Stand der Forschung.}
\label{tab:forschungsableitung}

\footnotesize
\setlength{\tabcolsep}{4pt}
\renewcommand{\arraystretch}{1.08}

\begin{tabular}{@{}>{\raggedright\arraybackslash}p{3.7cm}>{\raggedright\arraybackslash}p{4.5cm}>{\raggedright\arraybackslash}p{5.0cm}@{}}
\toprule
\textbf{Forschungsbereich}
&
\textbf{Bezug zur vorliegenden Arbeit}
&
\textbf{Umsetzung}
\\
\midrule

Selbstüberwachte URL-Repräsentationen
(vgl. [Maneriker et al. (2021), PDF-S. 1--2]; [Wang et al. (2023), Abstract, o.\,S.]; [Li et al. (2025), S. 1--2 und 12])
&
Zusätzliche task-adaptive Anpassung bereits vortrainierter URL- und
HTML-Text-Repräsentationen.
&
Vergleich von R0, RU, RT und RUT bei konstanten Downstream-Probes.
\\[1mm]

Nutzung ungelabelter Daten und reduzierte Labelbudgets
(vgl. [Gui et al. (2024), S. 1]; [Li et al. (2025), S. 1--2 und 12])
&
Verhalten der Repräsentationsvarianten bei unterschiedlicher gelabelter
Downstream-Supervision.
&
Sieben Labelbudgets sowie Bestimmung der kleinsten ausreichenden
RUT-Stufe gegenüber vollständig mit 200.000 Labels trainierten Referenzen.
\\[1mm]

Mehrquellenbasierte Webseitenklassifikation
(vgl. [Aljofey et al. (2022), S. 1--2 und 13]; [Opara et al. (2024), S. 2]; [Yoon et al. (2024), S. 2 und 13--15])
&
Überführung der ausgewählten Repräsentationsbasis in ein multimodales
Klassifikationssystem.
&
Sequenzielle Erweiterung um HTML-Text, DOM und lernbare Gating-Fusion.
\\[1mm]

Low-FPR- und gestufte Verarbeitung
(vgl. [Maneriker et al. (2021), PDF-S. 1--2]; [Liu et al. (2021), Abstract, o.\,S.]; [ENISA (2020), S. 7]; [Reddi et al. (2020), S. 5--6])
&
Bewertung des Gesamtsystems bei festgelegtem Betriebspunkt und
unterschiedlicher Verarbeitungstiefe.
&
CAL-basierter Low-FPR-Betriebspunkt, URL-first-Kaskade sowie
Effizienz- und Eskalationsmetriken.
\\

\bottomrule
\end{tabular}
\end{table}
